\documentclass[../../main]{subfiles}

\begin{document}

\section{Homomorfismos de Estruturas Algébricas}
\label{sec:matematica:algebra:estruturas-algebricas:homomorfismos}

Nesta seção consideramos estruturas algébricas como modelos de uma assinatura
\[
  \Sigma = \{\sigma_i : n_i\}_{i \in I},
\]
e introduzimos aplicações que preservam sua estrutura.

\subsection{Definição geral de homomorfismo}

Sejam $(A, \Sigma)$ e $(B, \Sigma)$ duas estruturas algébricas sobre a mesma assinatura.

Uma função
\[
  f : A \to B
\]
é chamada de \textbf{homomorfismo} se preserva todas as operações da assinatura, isto é:
\[
  f(\sigma_i^A(a_1,\dots,a_{n_i})) =
  \sigma_i^B(f(a_1),\dots,f(a_{n_i})),
\]
para todo símbolo $\sigma_i \in \Sigma$ e todos $a_1,\dots,a_{n_i} \in A$.

\subsection{Interpretação estrutural}

Um homomorfismo é uma aplicação que preserva a interpretação das operações da assinatura.

Em outras palavras, ele respeita a forma como as operações são realizadas em cada estrutura.

\subsection{Homomorfismos de anéis}

Se $A$ e $B$ são anéis, isto é, estruturas sobre a assinatura
\[
  \Sigma = \{+, \cdot : 2\},
\]
uma função $f : A \to B$ é um homomorfismo de anéis se:

\begin{itemize}
  \item preserva a adição:
    \[
      f(a+b) = f(a) + f(b);
    \]

  \item preserva a multiplicação:
    \[
      f(a \cdot b) = f(a) \cdot f(b);
    \]

  \item (caso unitário) preserva o elemento neutro multiplicativo:
    \[
      f(1_A) = 1_B.
    \]
\end{itemize}

\subsection{Núcleo e imagem}

Se $f : A \to B$ é um homomorfismo de anéis, define-se:

\textbf{Núcleo:}
\[
  \ker(f) = \{a \in A \mid f(a) = 0\}.
\]

\textbf{Imagem:}
\[
  \mathrm{Im}(f) = \{f(a) \mid a \in A\}.
\]

\subsection{Propriedades estruturais}

\begin{itemize}
  \item O núcleo de um homomorfismo é um ideal de $A$;
  \item A imagem de um homomorfismo é um subanel de $B$.
\end{itemize}

\subsection{Ideia de compatibilidade estrutural}

Homomorfismos são aplicações que respeitam a assinatura da estrutura, preservando todas as operações definidas em $\Sigma$.

Assim, eles permitem comparar estruturas algébricas de forma coerente.

\subsection{Observação conceitual}

A noção de homomorfismo unifica diversas construções algébricas:

\begin{itemize}
  \item em grupos, preserva a operação binária;
  \item em anéis, preserva adição e multiplicação;
  \item em corpos, preserva ambas as operações com inversos.
\end{itemize}

Essa generalidade será essencial para a formulação dos teoremas de isomorfismo e para a estruturação categórica da álgebra.

\end{document}
